\section{The CAGE-S Scheduler}
\label{sec:method}

\subsection{Method overview}

CAGE-S is an online scheduling framework for sensitive microservices
executing across heterogeneous edge, fog, and cloud nodes. Its design is
based on the observation that compliance, operational risk, placement
fairness, and resource executability are fundamentally different concerns.
They should therefore not be represented as interchangeable terms in a
single weighted objective.

CAGE-S organizes the scheduling decision as an ordered sequence of
admission and selection stages. The policy gate first determines which
nodes are allowed to host the service. The temporal-risk model then
estimates whether the currently admissible nodes are likely to cross the
operational utilization boundary during the next control interval. A
fairness-aware candidate reduction step limits repeated concentration on
the same destinations. Finally, a calibrated Q95 guardrail verifies the
projected CPU and memory headroom before the selected node is committed.

The resulting decision path is

\begin{equation}
\label{eq:cage-decision-path}
\text{policy gate}
\rightarrow
\text{temporal-risk estimation}
\rightarrow
\text{fair candidate reduction}
\rightarrow
\text{Q95 verification}
\rightarrow
\text{reservation and commitment}.
\end{equation}

This ordered structure is central to the method. Mandatory service
requirements define the admissible action space before any predictive
ranking is applied. Temporal risk is then used to distinguish stable
candidates from nodes that are approaching an operational boundary.
Selection history controls repeated destination use, while the final
guardrail verifies that the chosen destination retains sufficient
projected headroom at commitment time.

CAGE-S always evaluates the current source before considering migration.
A source that satisfies all admission conditions is retained. Migration
is initiated only when the source violates at least one policy,
temporal-risk, or guardrail condition. This source-first rule avoids
unnecessary movement while applying the same admission requirements to
retained and newly selected placements.

\begin{figure*}[t]
    \centering
    \includegraphics[width=0.96\textwidth]
    {cage_stages_only.pdf}
    \caption{The ordered CAGE-S scheduling pipeline. A service event
    passes through policy gating, temporal-risk estimation,
    fairness-aware shortlisting, Q95 pre-commit verification, and
    reservation-aware commitment.}
    \label{fig:cage-stages}
\end{figure*}

\subsection{Policy-gated admissibility}

For each service requirement profile, CAGE-S constructs a Boolean policy
mask over the ordered node set. Let

\begin{equation}
\label{eq:policy-mask}
    \mathbf{l}_{s}
    =
    \left[
        L(s,n_1),
        L(s,n_2),
        \ldots,
        L(s,n_M)
    \right]
    \in\{0,1\}^{M}
\end{equation}

denote the mask associated with service $s$, where the legality predicate
$L(s,n)$ is defined in Eq.~\eqref{eq:legal}. The mask simultaneously
encodes jurisdiction, trust, encryption, and trusted-execution
requirements.

Policy masks are independent of operational telemetry. They can therefore
be reused across multiple scheduling events and recomputed only when a
service requirement or node capability changes. For an event
$e_{t,r}$, the mask lookup immediately partitions the infrastructure into
admissible and prohibited destinations.

The scheduler first evaluates the current source node $n_{t,r}$. The
source can be retained only when

\begin{equation}
\label{eq:source-admission-method}
\begin{split}
    \operatorname{SourceOK}(e_{t,r})
    ={}&
    L(s_{t,r},n_{t,r})\\
    &{}\land
    \left(p_{n_{t,r},t}\leq\theta\right)\\
    &{}\land
    \left(g_{n_{t,r},t}\leq1\right).
\end{split}
\end{equation}

When this condition holds, CAGE-S preserves the current placement. When
it does not hold, the scheduler constructs the policy-admissible migration
set

\begin{equation}
\label{eq:base-candidate-method}
\begin{split}
    \mathcal{B}_{t,r}
    =
    \big\{
        n\in\N:
        &\;L(s_{t,r},n)=1,\\
        &\;n\notin A_t,\\
        &\;U_{t,r}(n)<C
    \big\}.
\end{split}
\end{equation}

The active-source condition $n\notin A_t$ prevents a destination from
coinciding with the source node of another event in the same control
batch. The reservation condition $U_{t,r}(n)<C$ ensures that simultaneous
decisions do not create conflicting commitments.

Only nodes in $\mathcal{B}_{t,r}$ proceed to the predictive and ranking
stages. Consequently, prohibited nodes are removed before temporal
inference influences the final decision.

\subsection{Temporal-risk estimation}

After policy-admissible candidates have been identified, CAGE-S estimates
their probability of crossing the operational utilization boundary during
the next control interval. The risk estimator is a histogram
gradient-boosting classifier

\begin{equation}
\label{eq:method-risk-model}
    p_{n,t}
    =
    f_{\boldsymbol{\phi}}
    \left(
        \mathbf{h}_{n,t}
    \right),
\end{equation}

where $\boldsymbol{\phi}$ denotes the learned model parameters and
$\mathbf{h}_{n,t}\in\R^{22}$ represents the recent operational history of
node $n$.

The temporal representation combines five complementary types of
information. Current-state features describe the observed CPU and memory
levels at decision time. Lagged features retain the two preceding
resource observations. Difference features encode recent direction and
acceleration. Short-window statistics summarize local level, range, and
stability. Projected-headroom features relate the current state to the
operational boundary.

The representation can be written as

\begin{equation}
\label{eq:temporal-feature-vector}
\begin{split}
    \mathbf{h}_{n,t}
    =
    \big[
        &\mathbf{x}_{n,t},
        \widehat{\mathbf{x}}^{\,\mathrm{up}}_{n,t+\Delta},
        u_{n,t},
        g_{n,t},\\
        &\mathbf{x}_{n,t-\Delta},
        \mathbf{x}_{n,t-2\Delta},\\
        &\Delta\mathbf{x}_{n,t},
        \Delta^{2}\mathbf{x}_{n,t},\\
        &\overline{\mathbf{x}}^{(3)}_{n,t},
        \operatorname{range}^{(3)}(\mathbf{x}_{n,t}),\\
        &\operatorname{var}^{(2)}(\mathbf{x}_{n,t}),
        \mathbf{q}_{n,t}
    \big],
\end{split}
\end{equation}

where $u_{n,t}$ denotes normalized joint utilization,
$\mathbf{q}_{n,t}$ denotes CPU and memory headroom, and the remaining
terms represent the temporal statistics.

The temporal classifier produces one risk probability for each evaluated
node. A node is operationally admissible when

\begin{equation}
\label{eq:temporal-admission-method}
    p_{n,t}\leq\theta.
\end{equation}

CAGE-S performs feature construction and model inference at timestamp
level. Rather than invoking the predictor separately for every event,
the scheduler identifies the union of nodes required by the current event
batch, constructs their temporal features once, and performs batched
inference. The resulting risk vector is reused throughout all decisions
at timestamp $t$.

\begin{algorithm}[t]
\caption{Timestamp-level temporal-state preparation}
\label{alg:temporal-preparation}
\KwIn{Node snapshots $X_t$, $X_{t-\Delta}$, and
$X_{t-2\Delta}$; temporal model $f_{\boldsymbol{\phi}}$;
calibrated margins $q_c^{0.95}$ and $q_m^{0.95}$}
\KwOut{Risk probabilities $p_t$ and guardrail values $g_t$}

\ForEach{$n\in\N$}{
    Construct the temporal feature vector
    $\mathbf{h}_{n,t}$\;
    $p_{n,t}\leftarrow
    f_{\boldsymbol{\phi}}(\mathbf{h}_{n,t})$\;
    $\widehat c^{\,\mathrm{up}}_{n,t+\Delta}
    \leftarrow c_{n,t}+q_c^{0.95}$\;
    $\widehat m^{\,\mathrm{up}}_{n,t+\Delta}
    \leftarrow m_{n,t}+q_m^{0.95}$\;
    $g_{n,t}\leftarrow
    \max\left\{
    \widehat c^{\,\mathrm{up}}_{n,t+\Delta}/\rho_c,
    \widehat m^{\,\mathrm{up}}_{n,t+\Delta}/\rho_m
    \right\}$\;
}
\Return{$p_t,g_t$}\;
\end{algorithm}

\subsection{Fairness-aware candidate reduction}

A large edge--cloud infrastructure may contain thousands of admissible
nodes. Ranking every candidate for every event would increase decision
cost and would not prevent repeated selection of the same destinations.
CAGE-S therefore introduces a fairness-aware candidate reduction stage.

Let $H_{t,r}(n)$ denote the cumulative number of earlier decisions
assigned to node $n$ before event $e_{t,r}$. The count incorporates both
previous timestamps and preceding events from the current timestamp. It
therefore represents the placement pressure accumulated by each node.

The policy-admissible candidates are ordered by

\begin{equation}
\label{eq:fair-preorder}
    \left(
        H_{t,r}(n),
        g_{n,t},
        \operatorname{id}(n)
    \right).
\end{equation}

The first key favours nodes that have been selected less frequently. The
second favours greater projected resource headroom among nodes with equal
historical use. The fixed identifier guarantees deterministic ordering.

The preliminary shortlist is defined as

\begin{equation}
\label{eq:method-shortlist}
    K_{t,r}
    =
    \operatorname{First}_{k}
    \left(
        \operatorname{lexsort}_{n\in\mathcal{B}_{t,r}}
        \left[
            H_{t,r}(n),
            g_{n,t},
            \operatorname{id}(n)
        \right]
    \right),
\end{equation}

where $k$ denotes the maximum shortlist size. If
$|\mathcal{B}_{t,r}|<k$, all policy-admissible candidates are retained.

The shortlisted candidates are then screened using both temporal risk and
projected headroom:

\begin{equation}
\label{eq:accepted-shortlist-method}
    \mathcal{A}_{t,r}
    =
    \left\{
        n\in K_{t,r}:
        p_{n,t}\leq\theta
        \land
        g_{n,t}\leq1
    \right\}.
\end{equation}

If the initial shortlist contains no accepted candidate, CAGE-S extends
the same screening process to the remaining nodes in
$\mathcal{B}_{t,r}\setminus K_{t,r}$. The shortlist therefore accelerates
the common decision path while preserving the ability to identify any
acceptable node in the complete policy-admissible set.

\begin{algorithm}[t]
\caption{Fairness-aware candidate shortlisting}
\label{alg:fair-shortlisting}
\KwIn{Base candidates $\mathcal{B}_{t,r}$;
historical counts $H_{t,r}$; risk probabilities $p_t$;
guardrails $g_t$; threshold $\theta$; shortlist size $k$}
\KwOut{Accepted candidate set $\mathcal{A}_{t,r}$}

Order $\mathcal{B}_{t,r}$ by
$\left(H_{t,r}(n),g_{n,t},\operatorname{id}(n)\right)$\;
$K_{t,r}\leftarrow$ first
$\min(k,|\mathcal{B}_{t,r}|)$ candidates\;
$\mathcal{A}_{t,r}\leftarrow
\{n\in K_{t,r}:
p_{n,t}\leq\theta
\land g_{n,t}\leq1\}$\;

\If{$\mathcal{A}_{t,r}=\varnothing$}{
    $\mathcal{R}_{t,r}\leftarrow
    \mathcal{B}_{t,r}\setminus K_{t,r}$\;
    $\mathcal{A}_{t,r}\leftarrow
    \{n\in\mathcal{R}_{t,r}:
    p_{n,t}\leq\theta
    \land g_{n,t}\leq1\}$\;
}
\Return{$\mathcal{A}_{t,r}$}\;
\end{algorithm}

\subsection{Destination ranking}

After candidate screening, all nodes in $\mathcal{A}_{t,r}$ satisfy the
mandatory policy, temporal-risk, guardrail, active-source, and reservation
conditions. CAGE-S then selects the preferred destination through
lexicographic ranking:

\begin{equation}
\label{eq:method-destination-ranking}
    n^{*}_{t,r}
    =
    \operatorname*{lexmin}_{n\in\mathcal{A}_{t,r}}
    \left(
        H_{t,r}(n),
        p_{n,t},
        g_{n,t},
        \operatorname{id}(n)
    \right).
\end{equation}

Historical selection count remains the primary ranking criterion. Among
nodes with equal historical use, the scheduler prefers the lower
predicted temporal risk. Guardrail utilization then distinguishes nodes
with different projected resource headroom, and the fixed node identifier
resolves exact ties.

Because mandatory admission has already been enforced, this ranking does
not balance legal and illegal actions or executable and non-executable
actions. It compares only destinations that have passed every required
condition.

\subsection{Q95 pre-commit guardrail}

The temporal-risk model and the pre-commit guardrail perform complementary
functions. The temporal model estimates the probability of near-term
threshold crossing from recent operational dynamics. The guardrail
provides a direct projected-headroom verification immediately before the
destination is committed.

For node $n$, the CPU and memory upper projections are

\begin{align}
\label{eq:method-cpu-upper}
    \widehat c^{\,\mathrm{up}}_{n,t+\Delta}
    &=
    c_{n,t}+q_c^{0.95},\\
\label{eq:method-memory-upper}
    \widehat m^{\,\mathrm{up}}_{n,t+\Delta}
    &=
    m_{n,t}+q_m^{0.95}.
\end{align}

The normalized guardrail score is

\begin{equation}
\label{eq:method-guardrail}
    g_{n,t}
    =
    \max
    \left\{
        \frac{
        \widehat c^{\,\mathrm{up}}_{n,t+\Delta}
        }{\rho_c},
        \frac{
        \widehat m^{\,\mathrm{up}}_{n,t+\Delta}
        }{\rho_m}
    \right\}.
\end{equation}

The candidate passes the guardrail when

\begin{equation}
\label{eq:method-guardrail-pass}
    g_{n,t}\leq1.
\end{equation}

After the final destination has been selected, CAGE-S repeats the policy,
temporal, guardrail, active-source, and reservation checks. The decision
is committed only when the destination continues to satisfy all required
conditions.

\begin{algorithm}[t]
\caption{Pre-commit destination verification}
\label{alg:precommit-verification}
\KwIn{Selected destination $n^{*}$; service $s$;
risk $p_{n^{*},t}$; guardrail $g_{n^{*},t}$;
active-source set $A_t$; reservations $U_{t,r}$;
capacity $C$}
\KwOut{Commit decision}

$\textit{valid}\leftarrow
L(s,n^{*})=1$\;
$\textit{valid}\leftarrow
\textit{valid}\land
p_{n^{*},t}\leq\theta$\;
$\textit{valid}\leftarrow
\textit{valid}\land
g_{n^{*},t}\leq1$\;
$\textit{valid}\leftarrow
\textit{valid}\land
n^{*}\notin A_t$\;
$\textit{valid}\leftarrow
\textit{valid}\land
U_{t,r}(n^{*})<C$\;

\If{$\textit{valid}$}{
    $U_{t,r}(n^{*})\leftarrow U_{t,r}(n^{*})+1$\;
    Commit $n^{*}$\;
}
\end{algorithm}

\subsection{Reservation-aware batch scheduling}

Multiple events may arrive at the same control timestamp. CAGE-S processes
them in a deterministic order and maintains a destination reservation
counter for the batch.

At the beginning of timestamp $t$, the scheduler initializes

\begin{equation}
    U_t(n)=0,
    \qquad
    \forall n\in\N.
\end{equation}

Whenever a migrated event is committed to node $n^{*}$, the counter is
updated as

\begin{equation}
\label{eq:reservation-update-method}
    U_{t,r+1}(n^{*})
    =
    U_{t,r}(n^{*})+1.
\end{equation}

All subsequent events in the same batch observe the updated value. The
reservation mechanism therefore incorporates decisions made earlier in
the timestamp before later events construct their own candidate sets.

The historical count is updated in the same way:

\begin{equation}
\label{eq:history-update-method}
    H_{t,r+1}(n^{*})
    =
    H_{t,r}(n^{*})+1.
\end{equation}

Reservation counts control concurrent destination use within the current
timestamp, whereas historical counts regulate repeated selection across
the complete scheduling horizon.

\subsection{Complete scheduling algorithm}

Algorithm~\ref{alg:cage-complete} presents the complete CAGE-S procedure.
Temporal risks and guardrail values are prepared once for the current
timestamp. Each event then follows the same source-retention, candidate
construction, shortlisting, ranking, verification, and commitment path.

\begin{algorithm*}[t]
\caption{CAGE-S online scheduling at timestamp $t$}
\label{alg:cage-complete}
\KwIn{Ordered event batch $\Eset_t$; node snapshots
$X_t,X_{t-\Delta},X_{t-2\Delta}$; policy masks
$\{\mathbf{l}_s\}$; temporal model $f_{\boldsymbol{\phi}}$;
threshold $\theta$; shortlist size $k$; destination capacity $C$}
\KwOut{Placement decision $\pi(e)$ for each $e\in\Eset_t$}

$(p_t,g_t)\leftarrow$
Algorithm~\ref{alg:temporal-preparation}\;
$A_t\leftarrow
\{n_{t,r}:e_{t,r}\in\Eset_t\}$\;
Initialize $U_t(n)\leftarrow0$ for all $n\in\N$\;

\ForEach{$e_{t,r}=(i_{t,r},s_{t,r},n_{t,r},t)
\in\Eset_t$ in deterministic order}{

    $\textit{sourceOK}\leftarrow
    L(s_{t,r},n_{t,r})
    \land
    p_{n_{t,r},t}\leq\theta
    \land
    g_{n_{t,r},t}\leq1$\;

    \eIf{$\textit{sourceOK}$}{
        $n^{*}\leftarrow n_{t,r}$\;
    }{
        $\mathcal{B}_{t,r}\leftarrow
        \{n\in\N:
        L(s_{t,r},n)=1,\,
        n\notin A_t,\,
        U_t(n)<C\}$\;

        $\mathcal{A}_{t,r}\leftarrow$
        Algorithm~\ref{alg:fair-shortlisting}
        $(\mathcal{B}_{t,r},
        H_{t,r},
        p_t,
        g_t,
        \theta,
        k)$\;

        \If{$\mathcal{A}_{t,r}=\varnothing$}{
            $\pi(e_{t,r})\leftarrow\varnothing$\;
            \textbf{continue}\;
        }

        $n^{*}\leftarrow
        \operatorname*{lexmin}_{n\in\mathcal{A}_{t,r}}
        \left(
            H_{t,r}(n),
            p_{n,t},
            g_{n,t},
            \operatorname{id}(n)
        \right)$\;

        Verify $n^{*}$ using
        Algorithm~\ref{alg:precommit-verification}\;
    }

    $\pi(e_{t,r})\leftarrow n^{*}$\;
    $H_{t,r+1}(n^{*})
    \leftarrow H_{t,r}(n^{*})+1$\;
}
\Return{$\{\pi(e):e\in\Eset_t\}$}\;
\end{algorithm*}

\subsection{Computational complexity}

Let $M$ denote the number of infrastructure nodes, $B_t$ the number of
events at timestamp $t$, $k$ the shortlist size, and $R_t$ the number of
distinct nodes included in timestamp-level model inference. Temporal feature construction and risk inference require
$O(R_t)$ time for a fixed feature dimension and a fixed trained ensemble.
Policy gating is implemented as Boolean-array lookup. For each migrated
event, constructing the base candidate set requires $O(M)$ operations.
Fairness-aware top-$k$ selection can be performed in
$O(M\log k)$ time, followed by at most $O(k)$ operations for final
candidate ranking.

The overall timestamp-level complexity is therefore

\begin{equation}
\label{eq:method-complexity}
    O
    \left(
        R_t+
        B_tM\log k
    \right),
\end{equation}

where $k\ll M$. Policy masks require $O(PM)$ storage for $P$ distinct
service-policy profiles. Risk vectors, guardrail vectors, historical
counts, reservation counters, and timestamp-level node state require
$O(M)$ additional memory. The implementation reduces repeated work by reusing policy masks,
constructing temporal features once per timestamp, batching model
inference, and expanding beyond the initial shortlist only when required.
These mechanisms allow CAGE-S to preserve an explicit multi-stage
decision process while remaining suitable for large online candidate
sets.